En el articulo \textit{Prediction of Electricity Demand in Weakly Interconnected Power System Using an Ensemble Time Series model qith a Bayesian Optimization Approach}, pronostica la demanda electrica de corto plazo en el Sistema Electrico NAcional de Mexico, centrada en la region Peninsular, que esta integrada por Yucatan, Campeche y Quintana Roo.

En el articulo hace unso de modelos como; Sarima, Facebook Prophet.
Se enfoca en el estudio de 5 nodos electricos:
\begin{itemize}
    \item Nodo 41
    \item Nodo 42
    \item Nodo 43
    \item Nodo 44
    \item Nodo 45
\end{itemize}
Sarima presenta ventajas ya que la demanda tiene una estructura repetida, con estacionalidad clara.

Entre las limitaciones de esta investigacion utiliza solamente un periodo de 24 horas ademas de que su metrica principal fue el MAE.
\begin{table}[H]
    \centering
    \begin{tabular}{c|c|c|c}
    Nodo& Prophet & Sarima  & Ensamble BO \\
    41 & 17.52 & 15.72   & 9.60 \\
    42 & 51.86 & 15.99 & 28.18 \\
    43& 36.46 & 11.93 & 16.54 \\
    44& 9.19 & 12.44 & 4.45 \\
    45& 2.78 & 1.18 & 1.21 \\
    \end{tabular}
    \caption{Tabla comparativa MAE de los diferentes metodos del articulo }
    \label{tab:tablart1}
    {\raggedright \small \textit{Nota:} Datos tomados de \cite{PANALESPEREZ2025652}.\par}
\end{table}